\documentclass[tikz,border=4pt]{standalone}
\input{common}
\usetikzlibrary{arrows.meta,positioning,calc}
\begin{document}
\begin{tikzpicture}[
    box/.style={draw,minimum height=7mm,inner sep=0pt},
    ptr/.style={-{Stealth[length=2mm]},thick},
    circ/.style={ptr,accent,dashed}]
  % 结点 = prev | data | next
  \newcommand{\cell}[3]{%
    \node[box,minimum width=6mm] (#1l) at (#2,0) {};
    \node[box,minimum width=9mm,right=0pt of #1l] (#1d) {#3};
    \node[box,minimum width=6mm,right=0pt of #1d] (#1r) {};
    \fill (#1l.center) circle (1.3pt); \fill (#1r.center) circle (1.3pt);}
  \cell{h}{0}{\small 哨兵}
  \cell{a}{3.2}{$a_1$}
  \cell{b}{6.4}{$a_2$}
  \cell{c}{9.6}{$a_3$}
  \foreach \u/\v in {h/a,a/b,b/c} {
    \draw[ptr] ($(\u r.center)+(0,0.12)$) -- ($(\v l.west)+(0,0.12)$);
    \draw[ptr] ($(\v l.center)-(0,0.12)$) -- ($(\u r.east)-(0,0.12)$);}
  % 循环：尾结点 next 指回哨兵，哨兵 prev 指向尾结点
  \draw[circ] (cr.center) to[out=70,in=110,looseness=0.32] node[above,font=\small] {$\fld{a_3}{next}$} (hl.north);
  \draw[circ] (hl.center) to[out=-70,in=-110,looseness=0.32] node[below,font=\small] {$\fld{\text{哨兵}}{prev}$} (cr.south);
\end{tikzpicture}
\end{document}
